% Recetas de Cafés
\newpage
Hay muchas maneras de preparar el café, dependiendo del tipo de grano entre otras cosas se puede apreciar de diversas profundidades de sabores, acidez, aromas...

Sin embargo, supongo que cada persona tendrá su preferencia en esto último, y disfrutar según las típicas recetas tradicionales...

\begin{center}
\includegraphics[scale=0.25]{cover_cafe}
\end{center}

\newpage
% Café Irlandés
\begin{recipe}[source={Barista de IG}]{Café Irlandés}
	\introduction{
		El café irlandés es un café muy rico para disfrutar sobre todo en clímas fríos. La primera vez que lo probé fue en Concepción en el café Cantabria.
	}
	\ingredients{
		\unit[1]{oz} & de almibar de canela \\
		\unit[2]{oz} & de whisky \\
        \unit[3]{oz} & de café \\
        \unit[1.5]{oz} & de crema de leche líquida
	}
	\preparation{
		\begin{enumerate}
			\item Agregar agua hirviendo a una taza para que adquiera la temperatura, y luego de unos segundos, retirar.
			\item Agregar el almibar.
			\item Agregar el whisky.
			\item Agregar el café.
			\item En una coctelera, batir por 10 segundos la crema de leche.
			\item Con ayuda de la cuchara, agregar lentamente la crema de leche en la taza.
			\item Para finalizar, se puede empolvorear con canela (opcional).
		\end{enumerate}
	}
\end{recipe}